% !Mode:: "TeX:UTF-8"
% !TEX program  = xelatex
\documentclass[a4paper]{IEEEtran}
\IEEEoverridecommandlockouts
\usepackage{ctex}
\usepackage{cite}
\usepackage{amsmath,amssymb,amsfonts}
\usepackage{algorithmic}
\usepackage{graphicx}
\usepackage{textcomp}
\usepackage{xcolor}
\usepackage[colorlinks,linkcolor=blue]{hyperref}
\usepackage[a4paper, top=2.54cm, bottom=2.54cm, left=3.18cm, right=3.18cm]{geometry}
\title{\huge{现代演化理论与社会启发}\\
	\Large{——生物演化论与人类社会的发展}
}
\author{王凯灵 521030910356}
\begin{document}
	\maketitle
	\begin{abstract}
		生物演化理论的完整阐述最初来源于达尔文的《物种起源》。随人类社会发展，生物溯源技术进步，尤其是21世纪以来分子生物学的快速发展，生物演化理论不断修改完善，逐渐形成更加完整的理论体系。同时，生物演化论的内容也对人类社会产生深远影响，既有技术上的启发，也有宗教观点的矛盾。本文将从社会发展推动演化理论完善、演化理论与社会发展的关系两方面进行总结。
	\par\textbf{关键词—}
		生物演化、生物起源、LUCA假说、社会发展、遗传算法
	\end{abstract}

	\section{演化理论简述：\\社会条件与理论进步史}
		生物演化论的思想最早出现可追溯到公元前500年代的希腊哲学家Anaximander等人\cite{nationalgeo}，认为生物可以发生变异，并且随着时间的推移会发生适应性的改变。但是，由于西方世界长期受神学思想统治，主要的生物起源理论是神创论、进步论，意匠论等。落后的社会思想和宗教的绝对地位造成了生物起源研究长期被错误的学说支配。
	
		达尔文的著作《物种起源》发表于1859年。当此之时，欧洲第一次工业革命结束，多种科学思想爆发，神学理论在多个领域遭到挑战。此前，也有林奈、布丰，拉马克等人在博物学方面为进化理论打下基础\cite{origin}。
		
		同时，我认为地质学发展也对生物演化理论做出了极大贡献。由于地理大发现，地质学进入快速发展期。到了达尔文同时代，已经出现对于各个地质时代的基本划分，而地质研究中伴随的古生物化石，也就模糊的描绘出了一条动植物演化路线。
		
		近代，随着分子技术的发展，分子生物学可以将现存生物和古生物化石准确建立基因相似关系，刻画出了更加清晰的生物演化史。同时，我认为现代人类对于海洋的探索，地质的研究也十分重要，我从中梳理了一个完善的生物起源假说。
		
		生物演化论由古代哲学家和博物学家进行铺垫，在近代思想解放时期形成理论，由现代技术进步得到修正。生物演化理论的进步史也是人类社会发展史的一个重要侧面。
		
		\section{现代技术：\\社会发展与生物起源的新假说}
			在《地球生命》\cite{chujianjun}一书中提到的众多生物分类方法表明，生物分类方法一直以来饱受争议，其原因在于就可以确定的物种数目而言，传统的博物学发现的动植物和肉眼可见的菌落、藻类等微生物集体在整个演化树上只占微不足道的一小部分，但如果要将包括我们自己在内的动物与酵母菌等真核生物归类在一起，实在难以让常人接受。此外，含有真核生物蛋白基因的阿斯加德古菌的发现也一举挑战了所有的传统分类方法，让人困惑。生物演化论中，一个基本的假设是生命在地球上仅起源过一次，即所有生物都是同一个祖先演化的结果。因此，这个问题与生物如何起源的问题脱不开干系。
			
			在《地球生命》\cite{chujianjun}一书中，对生命起源的解释采取了原始汤的假说\cite{soup}，即生命起源于地表。然而，随着近几十年来的科技进步，我基于文献总结了一个更加完整的生命起源假说。
			\subsection{原始汤假说的致命缺陷}
			生命需要能量，而原始汤假说中，可以作为生命诞生原始能量的只有光能和闪电。其中，闪电含有的能量不稳定，不足以支持生命进行演化，而分子生物研究表明（如图\ref{fig:tree}），生物的最初自养形式是化能合成。至此，原始汤唯一的合理可能就是闪电辟出了一个可以化能自养的细胞，与其他假说相比，越发难以让人信服。	
			\begin{figure}[h]
				\centering
				\includegraphics[width=0.9\linewidth]{../tree}
				\caption{系统发生树[选自Wiki]，标记分支为化能古菌}
				\label{fig:tree}
			\end{figure}
			\subsection{LUCA海底起源假说}
			基于先前提出的共同祖先（LUCA）假说\cite{woesevolution, corliss}，
			结合文献\cite{hotbic,wave,smoke}和一些科普材料，我总结出了一条更优秀的生命起源路线。这一路线可以解释许多演化难题，如光合生物为何不采用能量最丰富的绿光。和《环太平洋》里的说法相似，对于生命起源，我们以前仰望天空，后来，我们看向海底。简单来说，早期海底火山喷涌出的矿物在迅速冷却下凝结时，会产生原始细胞大小的细孔。而海底火山的高温和海水的低温间的温度带中，存在一个合适生命诞生的区间。与闪电的效果类似，火山喷发时的能量也可以产生有机大分子，而矿物质细孔为细胞的产生提供了模具。细孔内外，酸碱中和反应不断提供稳定的能量，足够支持生命演化（具体而言，稳定能源和理化条件可以保证充足的物质基础，以及保证生命起源时产生的能量利用机制可以稳定延续，实现生命必须的逆熵增条件），而如今的三羧酸循环也正照应了当年利用矿物和酸碱中和反应的往事。
			\subsection{LUCA假说提出的社会背景}
			这一假说极大的依赖人类科技进步，尤其是21世纪以来。人们对于海洋的探测终于到达了深海海底，深度达到数千米。在这里，人们对深海生态展开研究\cite{hotbic,wave,smoke}，得到该假说的事实证据。分子生物学发展提供了生物溯源途径，确定了各种生命形式出现的大致顺序。地质学发展提供的地球历史大气含氧量曲线、地表矿物历史曲线也为LUCA假说和后续演化理论进行了重要补充。
			
			LUCA假说提出的时间是科学思想空前开放的现代，即使部分地区仍然难以接受演化理论，就整体而言，生物演化论的基本内容已经获得大众认可。
			
			\subsection{现代社会对生物演化论的补充}
			篇幅所限，我仅举光合作用的例子。在自然光谱中，绿色附近光谱含有能量最多（如图\ref{fig:sun}），但利用绿光进行光合作用的植物十分罕见，且多为光竞争中弱势的植物。经过理化分析，这一现象的直接原因是光合色素吸收绿光少、绿光波长的光量子的能量不容易激发形成ATP。但按照自然选择，在光合生物诞生之初，能吸收绿光的物种应该成为优势种存活，相关的能量机制也应演化成适合绿光的形式。这种与自然选择矛盾的普遍现象是生物演化论中的一大难题。
			\begin{figure}[h]
				\centering
				\includegraphics[width=0.9\linewidth]{../sun}
				\caption{自然光能量谱与5525K黑体辐射曲线[选自Nature]}
				\label{fig:sun}
			\end{figure}
		
			古生物研究发现，约30-40亿年前，地球表面大量存在的古菌含有的色素（如图\ref{fig:rel}）主要是紫色的类视黄醛物质\cite{pho}，即吸收了绿光，用于一些非光合的生物反应，产生酸碱性差异提供能量，与LUCA假说中的能量机制类似。此时，潜在的光合细菌被紫色菌毯覆盖，无法获得绿光。在菌毯之下，一些使用绿色素的细菌经过自然选择，产生了可以利用穿过菌毯的微弱光线光解水的色素，光解水产生的氢离子的作用取代了细菌所需的酸碱中和反应，同时给细菌带来充足的氢。然而，光合色素主要以有机结构卟啉为基础，而卟啉决定了在光合色素附近，没有能稳定转化绿光的结构，即如果光和生物想要演化出利用绿色的后代，需要经历一些无法很好利用光的不良中间体，而这些中间体显然会被选择淘汰。同时，含有紫色素的菌毯由于可以获得充足的大气含氢物质，没有光解水的选择压，没能选择出光合反应的后代，并且在光合作用产生后的约十亿年后因大氧化事件大受打击。利用绿光的光合总用从此走到了死胡同。这样一来，生物演化论的重要问题得到了合理解释。
			\begin{figure}[h]
				\centering
				\includegraphics[width=0.9\linewidth]{../relation}
				\caption{光合作用诞生示意图\cite{pho}}
				\label{fig:rel}
			\end{figure}
		\section{矛盾与启发：\\演化论对人类社会发展的影响}
			\subsection{社会思想}
			生物演化论从一开始，就是对神学的挑战。随着十九世纪以来，越来越多的演化证据被发现，越来越多的学者接受演化理论，西方的思想解放不断攻击神学的权威，影响了西方社会对宗教与科学的看法。同时，生物进化论影响了进步主义思想的興起。演化理论表明，生命有一个“进步”的方向。社会作为一种高级的生命活动形式，也会朝着更高级的方向“进步”。这激发了19世纪的进步主义理想，进而促进了马克思主义等社会理论的发展与完善，为人类哲学做出极大贡献。
			\subsection{哲学思想}
			研究生物的具体演化历史，对于社会发展有重要启示。其中，我认为最有意思的一段古生物史是埃迪卡拉纪。在这一时段，地球结束漫长的雪球地球时期，有氧藻类迅速大量繁殖，有机物食物唾手可得，发生了阿瓦隆大爆发，诞生了大量埃迪卡拉动物。由于这一时期还没有占据支配地位的动物掠食者，藻类极其充足，在这种条件下演化出的埃迪卡拉动物仅仅演化出高效的进食和繁殖能力，由于几乎不存在选择压，没有证据表明其曾演化出运动能力，感知能力，以及很可能此前已经在真后生动物中出现的神经系统。从某种意义上，这与理想的共产主义十分类似：物质基础十分富裕，没有社会竞争，社会产品按需分配。埃迪卡拉动物的结局是不复存在，也意味着当生产力水平与群众觉悟没有到达一定高度，如果社会强行进入社会主义和共产主义，只会带来社会财富的耗尽和新的矛盾与压迫，这对社会学研究具有重要启发意义。然而，当人类掌握人工智能，创造无欲无求，能够自发获取能量进行生产的机器，或许就有了在小小地球上无限逆转熵增，进入后共产主义社会的可能。
			\subsection{科技启发}
			演化论为人类社会的科学发展也做出了贡献。就我熟悉的计算机和人工智能领域而言，演化理论在算法上有两个十分直接的启发。
			
			其一是遗传算法。顾名思义，当我们想要寻找一个问题的最优解时，我们可以只给出解的比较方法，模拟自然选择过程：选取网格或边界值作为初始生物，按给定比较方法相互比较。此后每一轮，筛去不够优秀的亲本，将存活的亲本进行杂交，并设计重组和突变得到下一代，直到稳定，得到的解往往是全局最优解。
			
			其二是强化学习，对于一个人工智能模型，使其与一个环境交互，初始时模型的行为随机，按照一定条件给模型奖励。这就像自然选择奖励适合自然的生物活下来，不适合的生物死亡。这种方法得到的人工智能往往具有高度适应环境的能力，典型代表就是下围棋的AlphaGo模型。
		\section{总结}
			演化理论的思想自古就有，伴随人类社会发展、思想解放走向成熟。步入现代，许多学科的快速发展完善了演化理论，依据演化理论溯源生物，给出了LUCA、叶绿素等合理假说。演化理论也促进了思想解放，哲学和科技发展。生物演化论与人类社会发展有千丝万缕的联系，在人类史上至关重要。
	\begin{thebibliography}{99}
		\bibitem{nationalgeo} National Geographic Education, Theory of Evolution
		, National Geographic Society[O]. Available: \url{https://education.nationalgeographic.org/resource/theory-evolution/} [Accessed May 6, 2023].
		\bibitem{origin} 卢继传. 达尔文主义是怎样发展与修正的[J]. 生物学杂志, 1987, 5.
		\bibitem{chujianjun} 褚建君. 地球生命[M]. 上海交通大学出版社, 2015: 2-3.
		\bibitem{soup} Miller S L. A production of amino acids under possible primitive earth conditions[J]. Science, 1953, 117(3046): 528-529.
		\bibitem{corliss} Corliss J B, Baross J A, Hoffman S E. An hypothesis concerning the relationships between submarine hot springs and the origin of life on earth[J]. Oceanologica Acta, Special issue, 1981.
		\bibitem{woesevolution} Corliss J B, Dymond J, Gordon L I, et al. Submarine thermal springs on the Galapagos Rift[J]. Science, 1979, 203(4385): 1073-1083.
		\bibitem{hotbic} Canganella F, Wiegel J. Anaerobic thermophiles[J]. Life, 2014, 4(1): 77-104.
		\bibitem{wave} Ding Y, Batista B, Steinbock O, et al. Wavy membranes and the growth rate of a planar chemical garden: Enhanced diffusion and bioenergetics[J]. Proceedings of the National Academy of Sciences, 2016, 113(33): 9182-9186.
		\bibitem{smoke} Cardoso S S S, Cartwright J H E. On the differing growth mechanisms of black-smoker and Lost City-type hydrothermal vents[J]. Proceedings of the Royal Society A: Mathematical, Physical and Engineering Sciences, 2017, 473(2205): 20170387.
		\bibitem{pho} Blankenship R E. Early evolution of photosynthesis[J]. Plant physiology, 2010, 154(2): 434-438.
		\bibitem{funstuff} 芳斯塔芙. 鬼谷闲谈[O] \url{https://space.bilibili.com/72270557} [Accessed May 8, 2023]
	\end{thebibliography}
\end{document}